% =============================================================
%  Page geometry
% =============================================================
\usepackage[
    paperwidth  = 7.5in,
    paperheight = 9.25in,
    top         = 1.0in,
    bottom      = 1.1in,
    inner       = 1.1in,
    outer       = 1.0in,
    marginpar   = 0.7in,
    marginparsep= 0.15in
]{geometry}

% =============================================================
%  Typography
% =============================================================
\usepackage{lmodern}
\usepackage[T1]{fontenc}
\usepackage[utf8]{inputenc}
\usepackage{microtype}

% =============================================================
%  Colors  (load BEFORE tcolorbox and hyperref)
% =============================================================
\usepackage[dvipsnames]{xcolor}

\definecolor{TheoremBg}{RGB}{230, 240, 252}
\definecolor{TheoremFrame}{RGB}{0, 90, 160}
\definecolor{DefGreen}{RGB}{0, 130, 80}
\definecolor{DefBg}{RGB}{235, 252, 240}
\definecolor{LemmaYellow}{RGB}{255, 248, 220}
\definecolor{LemmaFrame}{RGB}{180, 140, 0}
\definecolor{PropBg}{RGB}{245, 235, 255}
\definecolor{PropFrame}{RGB}{120, 60, 180}
\definecolor{ExBg}{RGB}{252, 248, 240}
\definecolor{ExFrame}{RGB}{200, 140, 60}
\definecolor{RemBg}{RGB}{250, 245, 245}
\definecolor{RemFrame}{RGB}{180, 80, 80}
\definecolor{DarkNavy}{RGB}{10, 20, 60}
\definecolor{OffWhite}{RGB}{245, 245, 240}
\definecolor{Blue}{RGB}{0, 90, 160}
\definecolor{AccentGray}{RGB}{120, 130, 150}

\definecolor{PureWhite}{HTML}{FFFFFF}
\definecolor{PureBlack}{HTML}{000000}

\definecolor{OffWhite}{HTML}{F7F5EF}
\definecolor{DarkNavy}{HTML}{1F2A44}

\definecolor{WarmOffWhite}{HTML}{F4F1E8}
\definecolor{DeepNavy}{HTML}{1C274C}

\definecolor{DarkBg}{HTML}{0F172A}
\definecolor{LightText}{HTML}{F8FAFC}

\definecolor{SoftDarkBg}{HTML}{1E293B}
\definecolor{SoftLightText}{HTML}{E5E7EB}

% =============================================================
%  Math
% =============================================================
\usepackage{amsmath, amssymb, amsthm}
\usepackage{mathtools}
\usepackage{thmtools}

% =============================================================
%  Theorem boxes  (tcolorbox — load AFTER xcolor)
% =============================================================
\usepackage[most]{tcolorbox}

% ---- shared style helpers ----------------------------------
\tcbset{
    thmbase/.style = {
        enhanced, breakable,
        fonttitle = \bfseries\sffamily,
        separator sign = {\ $\triangleright$\ },
        description font = \normalfont\itshape,
        attach boxed title to top left = {yshift=-2mm, xshift=4mm},
        boxed title style = {
            sharp corners, size=small,
            colframe = #1, colback = #1
        },
        colframe = #1, sharp corners = south,
        top = 4pt, bottom = 4pt,
        before skip = 8pt, after skip = 8pt,
    }
}

% ---- Theorem -----------------------------------------------
\newtcbtheorem[number within=section]{theorem}{Theorem}{
    thmbase = TheoremFrame,
    colback  = TheoremBg,
    coltitle = white,
}{thm}

% ---- Definition --------------------------------------------
\newtcbtheorem[use counter from=theorem]{definition}{Definition}{
    thmbase  = DefGreen,
    colback  = DefBg,
    coltitle = white,
}{def}

% ---- Lemma -------------------------------------------------
\newtcbtheorem[use counter from=theorem]{lemma}{Lemma}{
    thmbase  = LemmaFrame,
    colback  = LemmaYellow,
    coltitle = white,
}{lem}

% ---- Proposition -------------------------------------------
\newtcbtheorem[use counter from=theorem]{proposition}{Proposition}{
    thmbase  = PropFrame,
    colback  = PropBg,
    coltitle = white,
}{prop}

% ---- Corollary ---------------------------------------------
\newtcbtheorem[use counter from=theorem]{corollary}{Corollary}{
    thmbase  = TheoremFrame,
    colback  = TheoremBg!50!white,
    coltitle = white,
}{cor}

% ---- Example -----------------------------------------------
\newtcbtheorem[use counter from=theorem]{example}{Example}{
    thmbase  = ExFrame,
    colback  = ExBg,
    coltitle = white,
}{ex}

% ---- Remark ------------------------------------------------
\newtcbtheorem[use counter from=theorem]{remark}{Remark}{
    thmbase  = RemFrame,
    colback  = RemBg,
    coltitle = white,
}{rem}

% ---- Proof -------------------------------------------------
\tcolorboxenvironment{proof}{
    blanker, breakable,
    left = 6mm,
    before skip = 4pt, after skip = 8pt,
    borderline west = {1.5pt}{0pt}{AccentGray},
}

% =============================================================
%  Section / Subsection formatting  (titlesec)
% =============================================================
\usepackage{titlesec}
\usepackage{titletoc}

\titleformat{\section}
    {\large\bfseries\sffamily\color{DarkNavy}}
    {\color{Blue}\thesection}{0.8em}{}
    [{\color{Blue}\titlerule[0.6pt]}]

\titleformat{\subsection}
    {\normalsize\bfseries\sffamily\color{DarkNavy}}
    {\color{Blue}\thesubsection}{0.7em}{}

\titleformat{\subsubsection}[runin]
    {\normalsize\bfseries\sffamily\color{AccentGray}}
    {\thesubsubsection}{0.6em}{}[.\quad]

\titlespacing*{\section}      {0pt}{1.8ex plus .4ex}{0.8ex plus .2ex}
\titlespacing*{\subsection}   {0pt}{1.2ex plus .3ex}{0.5ex plus .1ex}
\titlespacing*{\subsubsection}{0pt}{0.8ex plus .2ex}{0pt}

% =============================================================
%  Headers & Footers
% =============================================================
\usepackage{fancyhdr}
\pagestyle{fancy}
\fancyhf{}
\fancyhead[LE]{\small\sffamily\color{AccentGray}\thepage\ \textbar\ \leftmark}
\fancyhead[RO]{\small\sffamily\color{AccentGray}\rightmark\ \textbar\ \thepage}
\fancyfoot[C]{\small\sffamily\color{AccentGray}\textit{\docsubtitle}}
\renewcommand{\headrulewidth}{0.4pt}
\renewcommand{\footrulewidth}{0.3pt}
\renewcommand{\headrule}{\color{AccentGray}\hrule width\headwidth height\headrulewidth}
\renewcommand{\footrule}{\color{AccentGray}\hrule width\headwidth height\footrulewidth}

% Plain style for chapter/title pages
\fancypagestyle{plain}{
    \fancyhf{}
    \fancyfoot[C]{\small\sffamily\color{AccentGray}\thepage}
    \renewcommand{\headrulewidth}{0pt}
}

% =============================================================
%  Lists
% =============================================================
\usepackage{enumitem}
\setlist[itemize]{
    leftmargin=*, topsep=3pt, itemsep=2pt,
    label={\color{Blue}\textbullet}
}
\setlist[enumerate]{
    leftmargin=*, topsep=3pt, itemsep=2pt,
    label={\color{Blue}\arabic*.}
}

% =============================================================
%  Hyperlinks  (load near-last, before cleveref)
% =============================================================
\usepackage[
    colorlinks = true,
    linkcolor  = Blue,
    citecolor  = DefGreen,
    urlcolor   = Blue,
    pdfusetitle
]{hyperref}
\usepackage{cleveref}   % \cref{} — load AFTER hyperref

% =============================================================
%  Utilities
% =============================================================
\usepackage{graphicx}
\usepackage{booktabs}
\usepackage{array}
\usepackage{marginnote}
\renewcommand{\marginfont}{\footnotesize\sffamily\color{AccentGray}}
\usepackage{makeidx}
\makeindex

% =============================================================
%  Macros — sets & fields
% =============================================================
\newcommand{\R}{\mathbb{R}}
\newcommand{\C}{\mathbb{C}}
\newcommand{\N}{\mathbb{N}}
\newcommand{\Z}{\mathbb{Z}}
\newcommand{\Q}{\mathbb{Q}}
\newcommand{\F}{\mathbb{F}}

% --- Calculus / Analysis ------------------------------------
\newcommand{\dd}{\,\mathrm{d}}
\newcommand{\dmu}{\,\mathrm{d}\mu}
% FIX: \dx and \dt now accept an optional variable argument so that
%      \dx    => \,\mathrm{d}x   (default)
%      \dx[y] => \,\mathrm{d}y   (used e.g. in the convolution integral)
%      \dt    => \,\mathrm{d}t   (default)
\newcommand{\dx}[1][x]{\,\mathrm{d}#1}
\newcommand{\dt}[1][t]{\,\mathrm{d}#1}
\newcommand{\pdv}[2]{\dfrac{\partial #1}{\partial #2}}
\newcommand{\dv}[2]{\dfrac{\mathrm{d} #1}{\mathrm{d} #2}}

% --- Operators ----------------------------------------------
\DeclareMathOperator{\supp}{supp}
\DeclareMathOperator{\esssup}{ess\,sup}
\DeclareMathOperator{\essinf}{ess\,inf}
\DeclareMathOperator{\Leb}{Leb}
\DeclareMathOperator{\Dom}{Dom}
\DeclareMathOperator{\Ran}{Ran}
\DeclareMathOperator{\cl}{cl}
\DeclareMathOperator{\interior}{int}
\DeclareMathOperator{\sgn}{sgn}

% --- Paired delimiters (auto-scaling with *) ----------------
\DeclarePairedDelimiter{\norm}{\lVert}{\rVert}
\DeclarePairedDelimiter{\abs}{\lvert}{\rvert}
\DeclarePairedDelimiter{\inner}{\langle}{\rangle}
\DeclarePairedDelimiter{\ceil}{\lceil}{\rceil}
\DeclarePairedDelimiter{\floor}{\lfloor}{\rfloor}
\DeclarePairedDelimiter{\set}{\{}{\}}   % \set*{\,x : P(x)\,}

% --- Restriction --------------------------------------------
\newcommand{\restr}[2]{\left.\kern-\nulldelimiterspace #1\right|_{#2}}

% --- Misc ---------------------------------------------------
\newcommand{\eps}{\varepsilon}
\newcommand{\ph}{\varphi}
\newcommand{\wt}[1]{\widetilde{#1}}
\newcommand{\wh}[1]{\widehat{#1}}
\newcommand{\ol}[1]{\overline{#1}}

% Margin note shortcut
\newcommand{\mn}[1]{\marginnote{#1}}
